\section{Mật mã tính toán}

\subsection{(15 điểm) Các khái niệm về an toàn tính toán}

\textbf{(a) (5 điểm) Giải thích tại sao an toàn tính toán lại quan trọng trong thực tế mặc dù đã tồn tại khái niệm an toàn thông tin tuyệt đối (information-theoretic security). Thảo luận các hạn chế của cả hai phương pháp.}

\begin{center}
\textbf{Trả lời:}
\end{center}

\textbf{Tại sao Computational Security quan trọng:}
\begin{itemize}
    \item \textbf{Perfect Secrecy không thực tế}: OTP yêu cầu khóa dài bằng thông điệp và chỉ dùng một lần, bất khả thi cho Internet và hệ thống lớn.
    \item \textbf{Computational Security thực tế}: Cho phép khóa ngắn cố định (128-256 bit), có thể tái sử dụng, hiệu quả về mặt tính toán.
    \item \textbf{Cân bằng hợp lý}: Chấp nhận xác suất tấn công nhỏ (negligible) để đạt được tính khả thi thực tế.
\end{itemize}

\textbf{Hạn chế của Perfect Secrecy:}
\begin{itemize}
    \item Yêu cầu khóa dài bằng thông điệp $\rightarrow$ không thể triển khai quy mô lớn
    \item Không thể tái sử dụng khóa $\rightarrow$ vấn đề phân phối khóa phức tạp
    \item Chỉ đảm bảo confidentiality, không có authentication/integrity
\end{itemize}

\textbf{Hạn chế của Computational Security:}
\begin{itemize}
    \item Dựa vào giả định bài toán khó $\rightarrow$ có thể bị phá vỡ bởi công nghệ tương lai (ví dụ: máy tính lượng tử)
    \item Không đảm bảo an toàn tuyệt đối $\rightarrow$ vẫn có xác suất tấn công thành công (dù rất nhỏ)
    \item Phụ thuộc vào giới hạn tài nguyên tính toán hiện tại
\end{itemize}

\textbf{Kết luận}: Computational Security là giải pháp thực tế duy nhất cho mật mã học hiện đại, cho phép cân bằng giữa an toàn và tính khả thi.

\textbf{Tham khảo:} Chương 1.3 – Phần 1.2 (Mật mã học Tính toán), Chương 1.2 – Phần 4.4 (Vấn đề Phân phối Khóa).

\textbf{(b) (5 điểm) Xét tấn công vét cạn (brute-force) lên AES-128:}

\begin{itemize}
    \item \textbf{i.} Sử dụng bảng chi phí tiền tệ đã cung cấp trong bài giảng, ước tính chi phí để thử tất cả các khóa có thể.
    \item \textbf{ii.} Thảo luận liệu cách tiếp cận an toàn tính toán có hợp lý không khi xét đến chi phí này.
\end{itemize}

\begin{center}
\textbf{Trả lời:}
\end{center}

\textbf{i. Ước tính chi phí brute-force AES-128:}
\begin{itemize}
    \item Số khóa có thể: $2^{128} \approx 3.4 \times 10^{38}$
    \item Giả sử mỗi phép thử tốn $10^{-6}$ USD (rất lạc quan)
    \item Chi phí tổng: $2^{128} \times 10^{-6} \approx 3.4 \times 10^{32}$ USD
    \item So sánh: GDP toàn cầu năm 2023 $\approx$ $100$ nghìn tỷ USD = $10^{14}$ USD
    \item Tỷ lệ: Chi phí brute-force $\approx$ $3.4 \times 10^{18}$ lần GDP toàn cầu
\end{itemize}

\textbf{ii. Đánh giá tính hợp lý của Computational Security:}
\begin{itemize}
    \item \textbf{Rất hợp lý}: Chi phí brute-force vượt xa khả năng tài chính của toàn nhân loại
    \item \textbf{Margin an toàn lớn}: Ngay cả khi giảm chi phí xuống $10^{10}$ lần, vẫn không khả thi
    \item \textbf{So sánh với các mối đe dọa khác}: Chi phí tấn công side-channel, implementation bugs thấp hơn nhiều
    \item \textbf{Kết luận}: Computational Security với AES-128 cung cấp mức bảo vệ đủ mạnh cho hầu hết ứng dụng thực tế
\end{itemize}

\textbf{Tham khảo:} Chương 1.3 – Phần 2.2 (Thời gian khả thi và bất khả thi), Chương 1.3 – Phần 2.3 (Khả năng bỏ qua).

\textbf{(c) (5 điểm) Nghịch lý ngày sinh (``birthday paradox'') rất quan trọng để hiểu nhiều tấn công mật mã. Nếu một hàm băm cho ra đầu ra dài $n$ bit:}

\begin{itemize}
    \item \textbf{i.} Xấp xỉ cần bao nhiêu đầu vào ngẫu nhiên để băm thì mới có xác suất 50\% tìm được một cặp va chạm?
    \item \textbf{ii.} Hàm băm cần bao nhiêu bit đầu ra để đủ an toàn trước các tấn công kiểu birthday trong thập kỷ tới?
\end{itemize}

\begin{center}
\textbf{Trả lời:}
\end{center}

\textbf{i. Số đầu vào cần thiết cho xác suất 50\% collision:}
\begin{itemize}
    \item Công thức birthday paradox: Cần khoảng $\sqrt{N}$ mẫu để có xác suất 50\% collision
    \item Với $n$ bit đầu ra: $N = 2^n$ giá trị có thể
    \item Số đầu vào cần: $\sqrt{2^n} = 2^{n/2}$
    \item \textbf{Ví dụ cụ thể:}
    \begin{itemize}
        \item SHA-256 ($n=256$): Cần $2^{128}$ đầu vào $\approx$ $3.4 \times 10^{38}$ đầu vào
        \item SHA-1 ($n=160$): Cần $2^{80}$ đầu vào $\approx$ $1.2 \times 10^{24}$ đầu vào
    \end{itemize}
\end{itemize}

\textbf{ii. Kích thước đầu ra an toàn cho thập kỷ tới:}
\begin{itemize}
    \item \textbf{Mục tiêu}: Đảm bảo $2^{n/2}$ operations bất khả thi về mặt tính toán
    \item \textbf{Ước tính khả năng tính toán}: Giả sử $2^{80}$ operations là giới hạn thực tế trong 10 năm tới
    \item \textbf{Yêu cầu}: $2^{n/2} \geq 2^{80}$ $\rightarrow$ $n \geq 160$ bit
    \item \textbf{Khuyến nghị an toàn}: $n \geq 256$ bit (như SHA-256) để có margin an toàn
    \item \textbf{Lý do}: Công nghệ tính toán có thể tiến bộ nhanh hơn dự kiến
\end{itemize}

\textbf{Kết luận}: Hàm băm cần ít nhất 256 bit đầu ra để đảm bảo an toàn trước birthday attacks trong thập kỷ tới.

\textbf{Tham khảo:} Chương 1.3 – Phần 3.3 (Tấn công Sinh nhật), Chương 1.3 – Phần 2.2 (Thời gian bất khả thi).

\subsection{(15 điểm) Phân biệt}

\textbf{(a) (7 điểm) Xét hai thư viện sau:}

\textbf{L1:}
\begin{verbatim}
SAMPLE():
  X ← {0,1}^n
  Y ← X \oplus 1^n
  trả về (X, Y)
\end{verbatim}

\textbf{L2:}
\begin{verbatim}
SAMPLE():
  Y ← {0,1}^n
  X ← Y \oplus 1^n
  trả về (X, Y)
\end{verbatim}

Sử dụng kỹ thuật hybrid proof để chứng minh hai thư viện này là không thể phân biệt. Mô tả rõ từng thư viện trung gian.

\begin{center}
\textbf{Trả lời:}
\end{center}

\textbf{Kết luận}: Hai thư viện L1 và L2 \textbf{không thể phân biệt} (Adv = 0).

\textbf{Phân tích trực tiếp:}
\begin{itemize}
    \item \textbf{L1}: $X \leftarrow \{0,1\}^n$, $Y = X \oplus 1^n$ $\rightarrow$ $(X, Y)$ với $Y = X \oplus 1^n$
    \item \textbf{L2}: $Y \leftarrow \{0,1\}^n$, $X = Y \oplus 1^n$ $\rightarrow$ $(X, Y)$ với $X = Y \oplus 1^n$
    \item \textbf{Quan sát}: Trong cả hai trường hợp, $X \oplus Y = 1^n$ luôn đúng
    \item \textbf{Phân phối}: Cả hai đều sinh ra cặp $(X, Y)$ với $X$ đều trên $\{0,1\}^n$ và $Y = X \oplus 1^n$
\end{itemize}

\textbf{Hybrid Proof (chi tiết):}
\begin{itemize}
    \item \textbf{Hybrid 0}: L1 - $X \leftarrow \{0,1\}^n$, $Y = X \oplus 1^n$
    \item \textbf{Hybrid 1}: L2 - $Y \leftarrow \{0,1\}^n$, $X = Y \oplus 1^n$
    \item \textbf{Chứng minh}: Hybrid 0 và Hybrid 1 có cùng phân phối đầu ra
    \begin{itemize}
        \item Trong Hybrid 0: $X$ đều, $Y = X \oplus 1^n$ $\rightarrow$ $Y$ cũng đều (vì XOR với hằng số)
        \item Trong Hybrid 1: $Y$ đều, $X = Y \oplus 1^n$ $\rightarrow$ $X$ cũng đều
        \item Cả hai đều sinh ra cặp $(X, Y)$ với $X \oplus Y = 1^n$ và $X$ đều
    \end{itemize}
\end{itemize}

\textbf{Kết luận}: Không có distinguisher nào có thể phân biệt L1 và L2 vì chúng sinh ra cùng phân phối đầu ra.

\textbf{Tham khảo:} Chương 1.3 – Phần 3.1 (Computational Indistinguishability), Chương 1.3 – Phần 3.2 (Hybrid Proof Technique).

\textbf{(b) (8 điểm) Xét một giao thức mà Alice muốn gửi tin nhắn mã hóa cho Bob. Họ sử dụng sơ đồ sau:}

\textbf{L\_real:}
\begin{verbatim}
K ← {0,1}^n

ENCRYPT(M):
  R ← {0,1}^n
  C1 = R
  C2 = K \oplus R \oplus M
  trả về (C1, C2)
\end{verbatim}

\textbf{L\_ideal:}
\begin{verbatim}
ENCRYPT(M):
  C1 ← {0,1}^n
  C2 ← {0,1}^n
  trả về (C1, C2)
\end{verbatim}

\textbf{i.} Hai thư viện trên thực ra không thể phân biệt! Hãy xây dựng một bộ phân biệt.

\textbf{ii.} Đề xuất một chỉnh sửa tối thiểu để làm cho thư viện ``real'' ở trên trở nên an toàn và giải thích vì sao chỉnh sửa đó hiệu quả.

\begin{center}
\textbf{Trả lời:}
\end{center}

\textbf{i. Bộ phân biệt cho L\_real vs L\_ideal:}

\textbf{Distinguisher D:}
\begin{enumerate}
    \item Gọi $(C1, C2) \leftarrow \text{ENCRYPT}(0^n)$ (mã hóa chuỗi toàn 0)
    \item Kiểm tra: $C1 \oplus C2 = 0^n$
    \item Nếu đúng $\rightarrow$ xuất 1 (L\_real), ngược lại $\rightarrow$ xuất 0 (L\_ideal)
\end{enumerate}

\textbf{Phân tích:}
\begin{itemize}
    \item \textbf{Trong L\_real}: $C1 = R$, $C2 = K \oplus R \oplus 0^n = K \oplus R$
    \begin{itemize}
        \item $C1 \oplus C2 = R \oplus (K \oplus R) = K$ (khóa cố định)
        \item Xác suất $C1 \oplus C2 = 0^n$ = xác suất $K = 0^n$ = $2^{-n}$ (rất nhỏ)
    \end{itemize}
    \item \textbf{Trong L\_ideal}: $C1$ và $C2$ đều ngẫu nhiên độc lập
    \begin{itemize}
        \item $C1 \oplus C2$ là ngẫu nhiên đều trên $\{0,1\}^n$
        \item Xác suất $C1 \oplus C2 = 0^n$ = $2^{-n}$
    \end{itemize}
    \item \textbf{Lợi thế}: $\text{Adv} = |2^{-n} - 2^{-n}| = 0$ $\rightarrow$ \textbf{Không phân biệt được!}
\end{itemize}

\textbf{ii. Chỉnh sửa tối thiểu để làm L\_real an toàn:}

\textbf{Vấn đề}: $C1 = R$ làm lộ $R$, khiến $C2 = K \oplus R \oplus M$ có thể bị tấn công.

\textbf{Giải pháp}: Sử dụng hàm băm để che giấu $R$:
\begin{verbatim}
L_real_fixed:
  K ← {0,1}^n
  
  ENCRYPT(M):
    R ← {0,1}^n
    C1 = H(R)  // H là hàm băm
    C2 = K \oplus R \oplus M
    trả về (C1, C2)
\end{verbatim}

\textbf{Tại sao hiệu quả:}
\begin{itemize}
    \item \textbf{Che giấu R}: $H(R)$ không tiết lộ $R$ (tính một chiều của hàm băm)
    \item \textbf{Bảo vệ C2}: Không thể tính $K \oplus R$ từ $C1$ và $C2$
    \item \textbf{Tính đúng đắn}: Bob có thể giải mã bằng cách thử tất cả $R$ và kiểm tra $H(R) = C1$
    \item \textbf{An toàn}: Phân phối $(C1, C2)$ trở nên ngẫu nhiên và không thể phân biệt với L\_ideal
\end{itemize}

\textbf{Tham khảo:} Chương 1.3 – Phần 1.1 (Mô hình thư viện), Chương 1.3 – Phần 3.1 (Computational Indistinguishability), Chương 1.1 – Phần 2.2 (Hàm băm).